\section{财产犯罪 - 既遂和数额}
\begin{enumerate}
  \item 既遂。取得控制说。他人占有转为自己占有，建立自己的占有 \yjt 取得控制 \yjt 置于自己实际控制范围。
        \begin{enumerate}
          \item 小件物品。拿在身上。
          \item 大件物品。出上场门。
          \item \al{甲盗窃古董，扔到墙外，翻墙后寻不到}。
        \end{enumerate}
  \item 数额。
        \begin{enumerate}
          \item 分类。
                \begin{enumerate}
                  \item 特别巨大 \yjt\, 30W
                  \item 巨大 \yjt\, 3W
                  \item 较大 \yjt\, 2000 完整保护。
                  \item 数额不大。抢劫罪和特殊的盗窃。银行卡。
                  \item 很小 \yjt 治安处罚法。\al{多次偷一根葱、食堂抢同学鸡腿被打、偷外卖}。
                \end{enumerate}
          \item \zd{要主客观一致}。
                \begin{enumerate}
                  \item \al{天价葡萄案}。
                  \item \al{卖淫女将 3000 嫖资和大老板手表拿走}。
                  \item \al{甲在商场盗窃到女士乙的提包，发现里面只有100元}。\zd{还是有偷到较大数额的危害可能性}。
                  \item \al{甲欲盗窃某珠宝店 100 万元的珠宝，临时拿错，拿走旁边的价值 1 万元的珠宝}。
                \end{enumerate}
          \item 违禁品的数额。\al{抢到 30W 假币}。黑市价定。不好认定那就数额除以 10。
          \item 犯罪所得数额。不受事后状况影响。\al{偷 3 个煤气罐}。
        \end{enumerate}
  \item \zd{既遂数额以犯罪时为准，所得收益以犯罪后案发时为准}。
        \begin{enumerate}
          \item 不受事后主观价值评价影响。\al{甲抢夺乙手镯，裂开一半}。
          \item 不受事后处置行为影响。\al{甲盗的 10W 项链，只卖出去了 100}
        \end{enumerate}
  \item 犯罪成本不计入。\al{骗到 300，付出 100}。\zd{成本算犯罪工具}。
\end{enumerate}
